\scanpage{13}
du terme). On voit alors facilement que le caractère de $M$ coïncide, sur les éléments de $I$ \uline{d'ordre premier à $p$}, avec le \uline{caractère modulaire} de la représentation de $I$ \deleted{définie} sur $\unsure{\mathbb F_p|\mathbf F_p}$ définie par $M(\unsure{o|0})$\deleted{$\unsure{?}$}. Comme, pour $K$ de caractéristique nulle, la donnée d'un groupe de BT sur $K$ revient à la donnée d'une représentation $p$-adique de $\operatorname{Gal}(\overline K/K)$ (dans un Module libre de type fini sur $\underline{\mathrm Z}_p$), on peut se demander \deleted{si} s'il n'est pas vrai en tous cas que le caractère d'un $M$ (satisfaisant la condition de monodromie) est déterminé ainsi par la seule connaissance de $M(\unsure{o|0})={}_pM$. Il n'en est rien, car si cette condition était vérifiée, on en déduirait facilement, en utilisant la théorie de Brauer des relèvements virtuels des représentations sur $\underline{\mathrm F}_p$ en représentations sur \deleted{$\unsure{K}$} $\underline{\mathrm Z}_p$, que pour \uline{tout} $M_{\inserted{K}}$ satisfaisant la condition de monodromie, le caractère de $M$ coïncide avec le caractère modulaire de $M(\unsure{o|0})$ sur les éléments d'ordre premier à $p$ de $I$. En particu\deleted{\unsure{x}}lier, si $M$ a bonne réduction, il s'ensuivrait que le caractère modulaire \inserted{de $M(\unsure{o|0})$} serait trivial dès que $M$ a bonne réduction. Or c'est déjà faux en prenant $M=\underline{\mathrm Z}_p(1)$, donc $M(\unsure{o|0})=\mu_p$, \deleted{\unsure{sur}} en se plaçant dans le cas où $K$ ne contient pas les racines p.èmes de $1$. Dans le cas où il \deleted{\unsure{le}}\inserted{les} contient, au contraire, il n'y a pas d'ennuis pour les groupes de BT sur $K$ qui sur un $K'$ ont une bonne réduction qui est, soit de type constant (cf plus haut), soit de type multiplicatif, et il conviendrait alors d'examiner encore le cas où il est de type ind-unipotent sur la \deleted{fibre} point fermé.

\margin{\uline{NB} Cela impliquerait que la fonction caractère est à valeurs dans $\mathbb Z_p\subset W(k)$, dès que $K$ contient les racines p\unsure{ièmes} de $1$ [donc sur un \unsure{s-groupe} \unsure{convenable} de $I$, d'indice $\mid p-1$] \deleted{Cela \illegible{4}}}

Cas test: soient $E_i$\deleted{et}\inserted{$(i=1,2)$} deux courbes elliptiques sur $K$,qui ont bonne réduction essentielle, ${}_pE_1$ isomorphe à ${}_pE_2$, $K$ contient les racines p.èmes de $1$, est-il vrai que les fonctions traces $\phi_{E_1}$ et $\phi_{E_2}$ sur le groupe d'inertie (qu'on peut, comme on sait, calculer aussi via les $T_\ell$ pour $\ell$ nombre premier \uline{quelconque}) sont égales \inserted{\figure{\begin{tikzpicture}\node[draw,ellipse] {sur les éléments premiers à $p$};\end{tikzpicture}}}? Le \deleted{seul} cas intéressant est celui où l'une au moins des deux courbes, après extension $K\longrightarrow K'$, se réduit en une courbe à invariant de Hasse nul. Ou encore: si $A_{\inserted{K}}$ schéma abélien sur $K$
\mspage{12}
